\scanpage{86}
\mspage{17}
% archive foliation: 85
\uline{Formes positives}

Le thème étant géométrique, nous nous plaçons
sur un corps alg. clos $k$. Nous ne
considérons que des motifs $\frac12$ \uline{simples}
sur $k$, \inserted{purs} de poids $\rho$. Soit $M$ un tel
motif, considérons les formes
\[
M\times M\longrightarrow\mathbf Q(-\rho)
\]
i.e. les éléments de
\[
\Gamma_k\bigl(\mathcal Hom(M\otimes M,\mathbf Q(-\rho))\bigr)=\Phi(M)
\]
où le \deleted{$\Gamma_k$ \unsure{est}} $\mathcal Hom$ est bien de poids $0$.
Elles forment un vectoriel de \deleted{poids $0$}\inserted{dim finie}
sur $\mathbf Q$. On considère l'\inserted{\unsure{sous-}}espace des
formes \uline{symétriques} si
$\rho$ pair, \uline{alternées} si $\rho$ impair. Cet
espace sera noté $\Phi^s(M)$. Noter les
hom canoniques
\[
\Phi^s(M)\otimes\Phi^s(N)\longrightarrow\Phi^s(M\otimes N)
\]
